\scanpage{181}
\begin{center}- 30 -\end{center}\hfill IX

3.2.1. Notons tout de suite que la définition du modèle de Néron $A$
de $A_K$ est la même que dans le cas local, rappelée dans 1.1. L'existence
de ce modèle de Néron, et le fait qu'il est séparé de type fini sur $S$,
est bien connue, et se ramène en fait immédiatement au cas local.
Bien entendu, pour tout point \uline{fermé} $s\in S$, $S'=\operatorname{Spec}(O_{S,s})$ est un trait,
et $A\times_S S'$ n'est autre que le modèle de Néron de $A_K$ sur $S'$, dans le
sens envisagé dans 1.1.

Démontrons 3.2. Il est trivial que i) implique ii), car il
suffit de prendre $G=A$. Il reste donc à prouver que si $G$ est comme
dans (ii), alors $G\longrightarrow A$ est une immersion ouverte. Or c'est en effet
un cas particulier de 3.1 e).

\uline{Corollaire} 3.3. \uline{Avec les notations de} 3.2, \uline{supposons les composantes
neutres des fibres du modèle de Néron de rang unipotent nul}. \uline{Soit} $S'$
\uline{un schéma satisfaisant aux mêmes hypothèses que} $S$, $K'$ \uline{son corps des
fonctions rationnelles}, $f:S'\longrightarrow S$ \uline{un morphisme dominant},
$A_{K'}\simeq A_K\otimes_K K'\simeq$ \uline{fibre générique de} $A\times_S S'$, $A'$ \uline{le modèle de Néron de}
$A_{K'}$, \uline{et considérons le morphisme unique}
\[
 u:A\times_S S'\longrightarrow A'
\]
\uline{induisant l'identité sur les fibres génériques}. \uline{Alors} $u$ \uline{est une immersion
ouverte}, \uline{et induit donc un isomorphisme sur les composantes neutres}
\[
 (3.3.1)\qquad A^{\unsure{o|0}}\times_S S'\xrightarrow{\sim}A'^{\unsure{o|0}}\qquad.
\]

Il suffit en effet d'appliquer 3.2 dans le cas de $S'$, en
prenant $G\simeq A\times_S S'$.
\margin{348}
\margin{348}
\margin{180}
